\documentclass{article}
\usepackage[utf8]{inputenc}
\usepackage{graphicx}
\usepackage{caption}
\usepackage{amsmath}
\usepackage{amssymb}
\usepackage[a4paper, margin=2cm]{geometry}  % <--- Aggiunto per ridurre i bordi
\title{Algebra Lineare}
\author{Nicolò Giuliani}
\date{\today}

\begin{document}

\section{Generazione di uno spazio vettoriale}
Dei vettori $p_1,\dots,p_n$ generano uno spazio vettoriale V se:
\begin{equation*}
    A=\begin{pmatrix}
        p_1 \\
        \vdots \\
        p_n
    \end{pmatrix}
\end{equation*}
portata a scala con l'algoritmo di Gauss:
\begin{equation*}
    A'=\begin{pmatrix}
        q_1 \\
        \vdots \\
        q_z
    \end{pmatrix}
\end{equation*}
Quindi la base composta  dai vettori $\langle q_1,\dots,q_z\rangle$ hanno come dimensione $z$ e quindi generano tutti gli spazi vettoriali della stessa dimensione. \\
Per esempio $\langle q_1,\dots,q_z\rangle$ generano: \\
1) $\mathbb{R}^z$ \\
2) $\mathbb{R}_{z-1}[x]$ 
\vspace{2em}
\\Quindi dei vettori generano un sottospazio se la loro base ha la stessa dimensione del sottospazio.

\section{Trovare una base di uno spazio vettoriale partendo da vettori}
Sappiamo che $p_1,\dots,p_n \in W$. \\
Per prima cosa dobbiamo trovare una base V dei vettori, $p_1,\dots,p_n$:
\begin{equation*}
    A=\begin{pmatrix}
        p_1 \\
        \vdots \\
        p_n
    \end{pmatrix}
\end{equation*}
portandola a scala:
\begin{equation*}
    A'=\begin{pmatrix}
        q_1 \\
        \vdots \\
        q_z
    \end{pmatrix}
\end{equation*}
Quindi $\langle q_1,\dots,q_z\rangle$ genera V. \\
Quindi $dim(V)=z$. \\
Dato che i vettori appartengono a W, la loro base e' un sottospazio, di W, di dimensione $z$ se $dim(W)=z$; ovvero $V\subseteq W$. \\
Per GEL: $\langle q_1,\dots,q_z\rangle$ generano W.

\section{Calcolo della dimensione di $Ker \,F$}
Scriviamo la matrice A $\in M_{m,n}(\mathbb{R})$ associada ad F:
\begin{equation*}
    A=\begin{pmatrix}
        p_1 \\
        \vdots \\
        p_m
    \end{pmatrix}
\end{equation*}
semplifichiamola se serve:
\begin{equation*}
    A'=\begin{pmatrix}
        q_1 \\
        \vdots \\
        q_z
    \end{pmatrix}
\end{equation*}
calcoliamo il sistema ponendolo uguale ad zero:
\begin{align*}
    \begin{cases}
        q_1=0 \\
        \vdots \\
        q_z=0
    \end{cases}
\end{align*}
Quindi abbiamo $n$ incognite e $rr(A')=z$. Le soluzioni dipendono da $n-z$ parametri. \\
Quindi $dim(Ker \, F)=n-z$.

\section{Trovare una base di $Ker \, F$ e completarla ad una base di uno spazio vettoriale}
\subsection{Troviamo una base di $Ker \, F$}
Se A e' associata ad F:
\begin{align*}
    A=\begin{pmatrix}
        p_1 \\
        \vdots \\
        p_n
    \end{pmatrix}
\end{align*}
semplificandola:
\begin{align*}
    A'=\begin{pmatrix}
        q_1 \\
        \vdots \\
        q_z
    \end{pmatrix}
\end{align*}
Quindi $rr(A')=z$. \\
Calcoliamo il sistema uguale ad 0:
\begin{align*}
    \begin{cases}
        q_1=0 \\
        \vdots \\
        q_z=0
    \end{cases}
    \to 
    \begin{cases}
        x_1= \dots \\
        \vdots \\
        x_z= \dots
    \end{cases}
\end{align*}
Quindi $Ker \, F=\{(x_1,\dots,x_z) \mid x_i \in \mathbb{R}\}$. \\
Quindi scegliamo un valore di $x_i$ e lo sostituiamo: $Ker \, F=\{(x_1,\dots,x_z)\}$
\subsection{Completiamola ad una base di uno spazio vettoriale}
Completiamola ad $\mathbb{R}^y$: \\
\begin{align*}
    \mathbb{R}^y=\{(\vec{v_1}),\dots,(\vec{v_y})\}
\end{align*}
quindi prendiamo $Ker \, F=\vec{v_1}$:
\begin{align*}
    \mathbb{R}^y=\{(Ker \, F),\dots,(\vec{v_y})\}
\end{align*}
Scegliamo gli altri vettori che mancano:
\begin{align*}
    \vec{v_2}=(0_1,1,\dots,0_z) \\
    \vdots \\
    \vec{v_y}=(0_1,\dots,1_y,\dots,0_z)
\end{align*}
Quindi ora scriviamo la base con i vettori scelti:
\begin{align*}
    \beta=\{(Ker \, F),(\vec{v_2}),\dots,(\vec{v_y})\} \\
    =\{(Ker \, F),(0_1,1,\dots,0_z),\dots,(0_1,\dots,1_y,\dots,0_z)\}
\end{align*}
\section{Stabilire se $\vec{v}\in Im \, F$}
Scriviamo la matrice A e semplifichiamola:
\begin{align*}
    A'=\begin{pmatrix}
        q_1 \\
        \vdots \\
        q_z
    \end{pmatrix}
\end{align*}
scriviamo la matrice pari ad $\vec{v}=\{v_1,\dots,v_z\}$
\begin{align*}
    A'\mid\vec{v}=\begin{pmatrix}
        q_1 & \mid & v_1\\
        \vdots & \mid \\
        q_z & \mid & v_z
    \end{pmatrix}
\end{align*}
Se $rr(A')=rr(A'\mid\vec{v})\to$ il sistema ha soluzione $\to v \in Im \, F$. \\
Se $rr(A')\neq rr(A\mid\vec{v}) \to$ nessuna soluzione $\to v \notin Im \, F$. \\
\subsection{Calcolare $F^{-1}(v)$}
Scriviamo:
\begin{align*}
    A'\mid\vec{v}=\begin{pmatrix}
        q_1 & \mid & v_1\\
        \vdots & \mid \\
        q_z & \mid & v_z
    \end{pmatrix}
\end{align*}
\textit{Ricorda: }$rr(A')=z$ \\
Le soluzioni dipendono da $n-z$ parametri.
\begin{align*}
    \begin{cases}
        q_1=v_1 \\
        \vdots \\
        q_z=v_z
    \end{cases}
    \to 
    \begin{cases}
        x_1 = \dots \\
        \vdots \\
        x_z = \dots
    \end{cases}
\end{align*}
Quindi $F^{-1}(\vec{v})=\{(x_1,\dots,x_z) \mid x_i \in \mathbb{R}\}$

\section{Calcolare le cordinate di un vettore rispetto ad una base}
Calcoliamo le cordinate di $\vec{v}$ rispetto ad $\beta$:
\begin{align*}
    \vec{v}=\lambda_1(x_1)+\dots+\lambda_n(x_n)
\end{align*}
\textit{Ricorda: } $\beta=\{(x_1),\dots,(x_n)\}$ \\
\textit{Ricorda: }$x_i=\{x_{i1},x_{i2},\dots,x_{in}\}$. \\
Quindi calcoliamo il sistema:
\begin{align*}
    \begin{cases}
        \lambda_1x_{11}+\lambda_2x_{21}+\dots+\lambda_nx_{n1}=v_1 \\
        \vdots \\
        \lambda_1x_{1n}+\lambda_2x_{2n}+\dots+\lambda_nx_{nn}=v_n
    \end{cases}
    \to 
    \begin{cases}
        \lambda_1=\dots \\
        \vdots \\
        \lambda_n=\dots
    \end{cases}
\end{align*}
Le cordinate sono $(\lambda_1,\dots,\lambda_2)$.

\section{Trovare applicazione lineare $G:\mathbb{R}^3 \to \mathbb{R}^3$ tale che $T $◦$ G$}
Scriviamo $A\mid I$:
\begin{align*}
    A\mid I = \begin{pmatrix}
        v_1 & | & 1 &0&0\\
        v_2  &|&0&1&0\\
        v_3 &|&0&0&1 
    \end{pmatrix}\\
    \to \dots \\
    \to I \mid A^{-1}=\begin{pmatrix}
        1 &0&0 &|& p_1\\
        0&1&0 & |&p_2\\
        0&0&1 & | & p_3
    \end{pmatrix}
\end{align*}
Quindi $G=A^{-1}$.
\begin{align*}
    G=\begin{pmatrix}
        p_1\\p_2\\p_3
    \end{pmatrix}
\end{align*}

\section{Si determini la matrice $A_{\beta,C}$ }
Trovare la matrice $A_{\beta,C}$ rispetto al dominio $\beta=\{e_1-3e_2,4e_2-e_3,e_2\}$ nel dominio e la base canonica $C$ di $\mathbb{R}^3$ nel codominio.
\begin{align*}
 I_{\beta,C}=\begin{pmatrix}
    1&-3&0\\
    0&4&-1\\
    0&1&0
 \end{pmatrix}   
\end{align*}
Quindi:
\begin{align*}
    A_{\beta,C}=A*I_{\beta,C}\\
    =\begin{pmatrix}
        4&0&0\\
        0&5&1\\
        -3&4&1
    \end{pmatrix}\begin{pmatrix}
        1&-3&0\\
        0&4&-1\\
        0&1&0
     \end{pmatrix}  \\
     =\begin{pmatrix}
        4&0&0\\
        -15&21&5\\
        -15&15&4
     \end{pmatrix}
\end{align*}






\end{document}